% Options for packages loaded elsewhere
\PassOptionsToPackage{unicode}{hyperref}
\PassOptionsToPackage{hyphens}{url}
\PassOptionsToPackage{dvipsnames,svgnames,x11names}{xcolor}
\documentclass[
  10pt,
  fontset=fandol]{ctexart}
\usepackage{xcolor}
\usepackage[margin=16mm]{geometry}
\usepackage{amsmath,amssymb}
\setcounter{secnumdepth}{-\maxdimen} % remove section numbering
\usepackage{iftex}
\ifPDFTeX
  \usepackage[T1]{fontenc}
  \usepackage[utf8]{inputenc}
  \usepackage{textcomp} % provide euro and other symbols
\else % if luatex or xetex
  \usepackage{unicode-math} % this also loads fontspec
  \defaultfontfeatures{Scale=MatchLowercase}
  \defaultfontfeatures[\rmfamily]{Ligatures=TeX,Scale=1}
\fi
\usepackage{lmodern}
\ifPDFTeX\else
  % xetex/luatex font selection
\fi
% Use upquote if available, for straight quotes in verbatim environments
\IfFileExists{upquote.sty}{\usepackage{upquote}}{}
\IfFileExists{microtype.sty}{% use microtype if available
  \usepackage[]{microtype}
  \UseMicrotypeSet[protrusion]{basicmath} % disable protrusion for tt fonts
}{}
\makeatletter
\@ifundefined{KOMAClassName}{% if non-KOMA class
  \IfFileExists{parskip.sty}{%
    \usepackage{parskip}
  }{% else
    \setlength{\parindent}{0pt}
    \setlength{\parskip}{6pt plus 2pt minus 1pt}}
}{% if KOMA class
  \KOMAoptions{parskip=half}}
\makeatother
\setlength{\emergencystretch}{3em} % prevent overfull lines
\providecommand{\tightlist}{%
  \setlength{\itemsep}{0pt}\setlength{\parskip}{0pt}}
\usepackage{amsmath,amssymb,mathtools,needspace}
\xeCJKDeclareCharClass{CJK}{"2032,"0400->"04FF,"2160->"216B,"2460->"2473,"25A0,"25C7,"25CB,"25CF}
\IfFontExistsTF{Arial Unicode MS}{\xeCJKsetup{AutoFallBack=true}\setCJKfallbackfamilyfont{\CJKrmdefault}{Arial Unicode MS}}{}
\setcounter{secnumdepth}{0}
\setlength{\emergencystretch}{2em}
\allowdisplaybreaks
\usepackage{bookmark}
\IfFileExists{xurl.sty}{\usepackage{xurl}}{} % add URL line breaks if available
\urlstyle{same}
\hypersetup{
  pdftitle={后退的 Euler 格式},
  colorlinks=true,
  linkcolor={Maroon},
  filecolor={Maroon},
  citecolor={Blue},
  urlcolor={Blue},
  pdfcreator={LaTeX via pandoc}}

\title{后退的 Euler 格式}
\author{}
\date{}

\begin{document}
\maketitle

\subsubsection{5.2.2 后退的 Euler 格式}\label{ux540eux9000ux7684-euler-ux683cux5f0f}

方程 \(y'=f(x,y)\) 中含有导数项 \(y'(x)\)，这是微分方程的本质特征，也正是它难以求解的症结所在。数值解法的关键在于设法消除其导数项，这项手续称为\textbf{离散化}。由于差分是微分的近似运算，实现离散化的基本途径之一是直接用差商替代导数。譬如，若在点 \(x_n\) 列出方程

\[
y'(x_n)=f(x_n,y(x_n)),\tag{5.2.4}
\]

并用差商 \(\dfrac{y(x_{n+1})-y(x_n)}{h}\) 近似替代其中的导数 \(y'(x_n)\)，结果有

\[
y(x_{n+1})\approx y(x_n)+hf(x_n,y(x_n)).
\]

设 \(y(x_n)\) 的近似值 \(y_n\) 已知，用它代入上式右端进行计算，并取计算结果 \(y_{n+1}\) 作为 \(y(x_{n+1})\) 的近似值，这就是 Euler 格式。

对于在点 \(x_{n+1}\) 列出的方程 (5.1.1)，有

\[
y'(x_{n+1})=f(x_{n+1},y(x_{n+1})),
\]

若用向后差商 \(\dfrac{y(x_{n+1})-y(x_n)}{h}\) 替代导数 \(y'(x_{n+1})\)，则可将上式离散化得

\[
\frac{y_{n+1}-y_n}{h}=f(x_{n+1},y_{n+1}),
\]

即

\[
y_{n+1}=y_n+hf(x_{n+1},y_{n+1}),\tag{5.2.5}
\]

此为\textbf{后退的 Euler 格式}。

后退的 Euler 格式与 Euler 格式有着本质的区别，后者是关于 \(y_{n+1}\) 的一个直接的计算格式，这类格式是\textbf{显式的}；而格式 (5.2.5) 的右端含有未知的 \(y_{n+1}\)，它实际上是关于 \(y_{n+1}\) 的一个函数方程，这类格式是\textbf{隐式的}。

显式与隐式两类方法各有特点。考虑到数值稳定性等因素，人们有时需要选用隐式方法，但使用显式算法远比隐式方便。

隐式方程 (5.2.5) 通常用迭代法求解，而迭代过程的实质是逐步显式化。

设用 Euler 格式 \(y_{n+1}^{(0)}=y_n+hf(x_n,y_n)\) 给出迭代初值 \(y_{n+1}^{(0)}\)，用它代入式 (5.2.5) 的右端，使之转化为显式，直接计算得

\[
y_{n+1}^{(1)}=y_n+hf(x_{n+1},y_{n+1}^{(0)}),
\]

然后再用 \(y_{n+1}^{(1)}\) 代入式 (5.2.5) 的右端，又有

\href{../../source-images/pdf-122.jpeg}{原页}

\[
y_{n+1}^{(2)}=y_n+hf(x_{n+1},y_{n+1}^{(1)}).
\]

如此反复进行迭代，得

\[
y_{n+1}^{(k+1)}=y_n+hf(x_{n+1},y_{n+1}^{(k)})\qquad(k=0,1,\cdots).
\]

如果迭代过程收敛，则极限值 \(y_{n+1}=\lim\limits_{k\to\infty}y_{n+1}^{(k)}\) 必满足隐式方程 (5.2.5)，从而获得后退的 Euler 方法的解。

再考察后退的 Euler 格式的局部截断误差。假设 \(y_n=y(x_n)\)，则按式 (5.2.5) 有

\[
y_{n+1}=y(x_n)+hf(x_{n+1},y_{n+1}),\tag{5.2.6}
\]

由于

\[
f(x_{n+1},y_{n+1})=f(x_{n+1},y(x_{n+1}))+f_y(x_{n+1},\eta)[y_{n+1}-y(x_{n+1})],
\]

其中 \(\eta\) 介于 \(y_{n+1}\) 与 \(y(x_{n+1})\) 之间。又

\[
f(x_{n+1},y(x_{n+1}))=y'(x_{n+1})=y'(x_n)+hy''(x_n)+\cdots,
\]

代入式 (5.2.6) 有

\[
y_{n+1}=hf_y(x_{n+1},\eta)[y_{n+1}-y(x_{n+1})]+y(x_n)+hy'(x_n)+h^2y''(x_n)+\cdots,
\]

将它与 Taylor 展开式

\[
y(x_{n+1})=y(x_n)+hy'(x_n)+\frac{h^2}{2}y''(x_n)+\cdots
\]

相减，得

\[
y(x_{n+1})-y_{n+1}=hf_y(x_{n+1},\eta)[y(x_{n+1})-y_{n+1}]-\frac{h^2}{2}y''(x_n)+\cdots.
\]

再注意到

\[
\frac{1}{1-hf_y(x_{n+1},\eta)}=1+hf_y(x_{n+1},\eta)+\cdots,
\]

最后整理，得

\[
y(x_{n+1})-y_{n+1}\approx-\frac{h^2}{2}y''(x_n).\tag{5.2.7}
\]

\end{document}
